\clearpage
\VerdierPhysicalPageAnchor{0214}%
\begingroup\small\setlength{\parskip}{0pt}\setlength{\emergencystretch}{1.5em}
\setlength{\abovedisplayskip}{4pt plus 2pt minus 2pt}%
\setlength{\belowdisplayskip}{4pt plus 2pt minus 2pt}%

\VerdierSemanticAnchor{ch3d-3-2-5-b}
\noindent b) Dans le diagramme ci-après :

\VerdierSemanticAnchor{ch3d-3-2-5-4}
\begin{equation}
\resizebox{0.96\textwidth}{!}{%
$\begin{tikzcd}[ampersand replacement=\&,column sep=4.1em,row sep=3.0em]
\&\& Z''^\bullet
  \arrow[ddll,bend right=8,"{u(S^{m+n})}"']
  \arrow[dl,"\rho"']
  \arrow[dr,"\psi"]
  \arrow[ddrr,bend left=8,"{v(S^{m+n})}"]
\&\& \\
\& Z^\bullet
  \arrow[dl,"{u(S^n)}"']
  \arrow[dr,"{v(S^n)}"]
\&\& {Z'}^\bullet[n]
  \arrow[dl,"{u(S^m)[n]}"']
  \arrow[dr,"{v(S^m)[n]}"]
\& \\
X \&\& X'[n] \&\& X''[m+n]
\end{tikzcd}$%
}
\tag{3.2.5.4}
\end{equation}
le carré central et les triangles latéraux sont commutatifs dans
$\operatorname{comp}(\mathcal A)$.

\VerdierSemanticAnchor{ch3d-3-2-5-proof-conclusion}
En passant dans la catégorie $D(\mathcal A)$, le diagramme (3.2.5.4)
fournit l'égalité :
\[
\bigl(\delta^m(S^m(X',X''))\bigr)[n]\,
\delta^n(S^n(X,X'))
=\delta^{m+n}(S^{m+n}(X,X'')).
\]
Il suffit alors d'utiliser (3.2.3.2) pour achever la démonstration.

\VerdierSemanticAnchor{ch3d-3-2-6}
\medskip\noindent\textbf{3.2.6.}\quad
Un cas d'application particulièrement important de la proposition (3.2.5)
est le suivant. Soit :
\[
S^n(X,Y)=0\longrightarrow Y\xrightarrow{u_n}Z_{n-1}
\xrightarrow{u_{n-1}}\cdots\xrightarrow{u_1}Z_0
\xrightarrow{u_0}X\longrightarrow0
\]
un élément de $E_{\mathcal A}^{n}(X,Y)$. Posons pour tout entier $p$,
$1\leq p\leq n$ :

\VerdierSemanticAnchor{ch3d-3-2-6-1}
\begin{equation}
S_p^1=0\longrightarrow\operatorname{Ker}(u_{p-1})
\longrightarrow Z_{p-1}\longrightarrow\operatorname{Im}(u_{p-1})
\longrightarrow0.
\tag{3.2.6.1}
\end{equation}
On a alors l'égalité :

\VerdierSemanticAnchor{ch3d-3-2-6-2}
\begin{equation}
\delta^n(S^n(X,Y))=(-1)^{\frac{n(n-1)}2}
\delta^1(S_n^1)[n-1]\delta^1(S_{n-1}^1)[n-2]\cdots
\delta^1(S_1^1),
\tag{3.2.6.2}
\end{equation}
ainsi qu'on le démontre immédiatement en utilisant la proposition (3.2.5).

\VerdierSemanticAnchor{ch3d-3-2-7}
\medskip\noindent\textbf{3.2.7.}\quad
Buvons le calice jusqu'à la lie, et soient $X$ et $Y$ deux objets de
$\mathcal A$ et par conséquent de $\mathcal A^o$ : la catégorie opposée à
$\mathcal A$. Les ensembles $E_{\mathcal A}^{n}(X,Y)$ et
$E_{\mathcal A^o}^{n}(Y,X)$ sont égaux et on a deux applications :

\VerdierSemanticAnchor{ch3d-3-2-7-1}
\begin{equation}
\resizebox{0.74\textwidth}{!}{%
$\begin{tikzcd}[ampersand replacement=\&,column sep=5.0em,row sep=2.8em]
E_{\mathcal A}^{n}(X,Y)
  \arrow[r,"{\delta_{\mathcal A}^{n}}"]
  \arrow[d,equal]
\& \operatorname{Ext}_{\mathcal A}^{n}(X,Y) \\
E_{\mathcal A^o}^{n}(Y,X)
  \arrow[r,"{\delta_{\mathcal A^o}^{n}}"']
\& \operatorname{Ext}_{\mathcal A^o}^{n}(Y,X)
\end{tikzcd}$%
}
\tag{3.2.7.1}
\end{equation}
\endgroup
